%% Threats to validity (plan step 9.1). Written at the freeze, 28 October 2026, from
%% 00_Plan/Threats_and_Methods_Notes.md (notes 1 to 26) and the Phase 5 to 8 notes. Every number is
%% a macro from 05_Paper/tables/numbers.tex. Each paragraph names a choice, the risk it carries and
%% what we did about it.

Each methods choice below is stated with the risk it carries and what we did to bound it. We group them by the validity type they threaten.

\paragraph{Construct validity: what the ground truth means.} The labels come from an attack schedule (Section~\ref{sec:gt}). A benign conversation is one outside every attack window, not one that was checked to be benign, and the schedule's drop rules remove Thursday afternoon, so \cicids{} has no infiltration class here. The window ends we reuse differ from those of the NFS-2023 labeller (SSH-Patator ends \NsshWindowEndMinutes{} minutes earlier in our schedule), which moves a few conversations at window edges from attack to benign. A schedule label also makes every record of an attack conversation an attack record, including records an exporter writes after the attack payload has ended (Section~\ref{sec:gt}). We chose the schedule because it is the only label that is the same in every arm by construction, so the comparison isolates the exporter; the cost is that a detector is asked to flag a record that carries no attack payload. That cost is measurable, and it is not small. Under payload-aware labels the RQ3 drop disappears, the RQ4 residual on \cicids{} disappears, and H2 on \cicids{} no longer holds by the strict-majority rule (Sections~\ref{sec:rq2} to \ref{sec:rq4}). The exporter results on \cicids{} are therefore results about exporters under one label policy, the one that makes the arms comparable, and the label policy for records without payload is a second choice that an evaluator has to state. Each arm's own labels are kept as a further setting (Section~\ref{sec:rq4}); with them the releases agree far less on which model is best.

\paragraph{Construct validity: attempted attacks.} Conversations that the schedule calls an attack but that carry no payload stay under their schedule class in the primary analysis, because the attacker sent the traffic and an operator would want the attempt detected. CRiSIS-2022 and both NFS-2023 releases label most of these benign by policy (Section~\ref{sec:gt}). Folding them into benign is a frozen sensitivity setting. In RQ4 it changes the H4 verdict only when the folded labels are used for scoring a model trained without them; retrained on the folded labels, H4 is undecided as in the primary setting and H2 is unchanged, while \RQfourFoldRetrainedClasses{} of \RQfourClassesOf{} per-class verdicts at the conversation unit change. In RQ5 one reversal, WTMC-2021 on Web brute force, depends on this rule and is reported as fragile.

\paragraph{Construct validity: the unit.} A conversation is defined by the endpoint pair, the protocol and a silence of more than $T = \SetT$ seconds (Section~\ref{sec:key}). On \cicids{} every scored class has the same conversation count from \NsweepCicStableMin{} to \NsweepCicStableMax{} seconds; larger values merge distinct DDoS and DoS Hulk conversations because their tools reuse endpoint keys. On UNSW-NB15, Generic splits into \NsweepGenericFivePct{} percent more conversations at \NsweepCicStableMin{} seconds than at \SetT{}, so a shorter $T$ would change that class. The RQ4 conclusions are rerun over $T$ in Section~\ref{sec:rq4}. A conversation is also not independent of its neighbours: one attack campaign holds many conversations, and a campaign can straddle the split, as SSH-Patator does in every seed. Our grouped split removes leakage between records of one conversation, not between conversations of one campaign. The ``any'' rule turns a conversation into a false alarm when a single record is flagged, which penalises releases that write short tail records; the majority and byte-weighted rules are reported beside it.

\paragraph{Internal validity: the instrument.} The original \cicids{} release has a minute clock and does not take part in defining conversations; its records are mapped afterwards, and up to \NnondefAmbigMax{} of them per day have more than one candidate conversation. That arm also comes from the GeneratedLabelledFlows files, because the better known machine learning CSV files carry no five-tuple and no timestamp. The two have the same row count (\NcicRows{}), but we did not compare their feature values row by row, so a result on our CIC original arm transfers to the machine learning files only if their features agree. Clocks were reconciled by us against the published schedule (Section~\ref{sec:clocks}). A wrong offset would misalign whole windows; each offset was chosen so that the attack windows line up with the schedule, and the window edges that still differ after it are kept in the release clock table of the replication files rather than corrected. The record alignment was checked by hand on \NhandCases{} conversations chosen to cover every cardinality class; the check found \docnum{four} faults, all fixed and held by unit tests, but \NhandCases{} cases is a small sample, and the pairing rule could still be wrong on patterns the sample did not contain.

\paragraph{Internal validity: the releases.} Two choices about UNSW-NB15 change its numbers. Exact duplicate lines are dropped for RQ1 and RQ2 and kept for RQ5, where published results used them; near repeats are collapsed in a sensitivity arm only. The near repeat definition is ours (same five-tuple, start, end and byte count after exact duplicates are gone), and a different one would remove a different set of records; neither sensitivity arm changes the H2 verdict. The H2 verdict on UNSW-NB15 does depend on CIC-UNSW-NB15, the release whose records score lowest: without it, H2 does not hold on that benchmark (Section~\ref{sec:rq2}). The NFS-2023 pair differs downstream of the exporter as well, through a filter that the expiry rule triggers and through its own labeller, so RQ3 measures one generator setting together with the filter it triggers. Of the split behaviours traced in Section~\ref{sec:rq1}, the third FTP-Patator record in CIC original is only partly explained, and which benign records NFS-2023-TE writes that models take for DoS Hulk is not explained.

\paragraph{Internal validity: the detection experiment.} The grid runs on a subsample that keeps \SetBenignFraction{} of benign conversations and at most \SetAttackCap{} conversations per attack class. This raises attack prevalence, so precision and F1 are higher than they would be on the full data. The same conversations are kept in every arm, so the subsample does not favour one exporter, but absolute F1 values are not full-data values. Reweighting each test conversation to full-data prevalence leaves recall unchanged and lowers precision, on average over arms and models, by at most \RevHtPrecMeanMaxCic{} for a \cicids{} attack class and \RevHtPrecMeanMaxUnsw{} for a UNSW-NB15 one; H2 still holds on both benchmarks under these weights (Table~\ref{tab:rev-robust}). We also trained the random forest on all records of every arm, on seed \setnum{0} of the grouped split, and scored it on all test records. For \RevFullWithinFive{} percent of the \RevFullPairs{} pairs of arm and attack class with at least \setnum{20} test records, flow F1 changed by at most \setnum{0.05} (median change \RevFullFoneMedian{}), and precision fell by at most \RevFullPrecMaxFall{}. The exception is the UNSW-NB15 Fuzzers class, whose recall falls by \RevFullFuzzRecallMin{} to \RevFullFuzzRecallMax{} across arms when the model is trained with all benign traffic, so that its F1 falls by up to \RevFullFoneMaxAbs{}. Absolute UNSW-NB15 Fuzzers scores in this paper are therefore subsample scores; this check has one model and one seed and does not test the H2 verdict on the full data. The conversation-level contrasts of RQ5 are on the subsample, and we did not reweight them; the RQ5 table states which metrics depend on prevalence. Hyperparameters were tuned once per arm and model, on seed \setnum{0} of the grouped split, from a grid of \SetGridSize{} settings; a model that suits one release better under a wider search would show up here as part of the exporter or interaction share. LightGBM was given an L2 penalty after its training diverged on the first arm tuned, before any grid cell was run. Identifiers are removed from the features. Time-to-live fields in UNSW-NB15 could act as a shortcut; removing them from every UNSW-NB15 arm leaves the H2 verdict there unchanged (Section~\ref{sec:rq2}). The software is deterministic: \RQtwoReproCells{} cells rerun from scratch gave identical predictions, and the whole result set is regenerated from the raw releases in one run.

\paragraph{Internal validity: rebuilt published comparisons.} RQ5 rebuilds each claim as the authors framed it, but not always with their setup: we ran a subset of their algorithms where they tested several, used our own hyperparameters and feature lists where we did not take theirs, seeded models they left unseeded, and used NF-UNSW-NB15-v3 for the conversation-level check of a claim made on NF-UNSW-NB15-v2, which has no timestamps. Where our per-flow rebuild did not reproduce the published value, the verdict is about our rebuild of the comparison, not about the authors' run.

\paragraph{External validity.} Two benchmarks, both recorded in laboratory testbeds, and only published releases; no packet capture was re-extracted by us. Releases we could not obtain are not covered: the HERA versions of \cicids{} and every release of CSE-CIC-IDS2018. Exporters not represented among the releases, such as Zeek and Argus on \cicids, enter only through the literature. The models of the primary grid are three tabular learners, two of them tree ensembles, and the extension adds two more of a different kind; sequence models that read several records of a conversation at once may react differently to how an exporter cuts it.

\paragraph{Statistical conclusion validity.} The confirmatory family is H2 and H4, each judged once by a rule fixed before training; every other comparison, including the per-class verdicts that make up H2, is descriptive, and the per-class intervals are not adjusted for multiplicity. Bootstrap intervals resample conversations, so they treat conversations of one campaign as independent and are too narrow if anything; the alpha intervals share that limit. The variance shares of RQ2 are shares of the sum of squares, and the arm factor has six degrees of freedom on \cicids{} and two on UNSW-NB15 against two for the model. Omega squared, computed afterwards as a check, lies within \NomegaMaxGap{} of the share for every eligible small class and changes the direction of \NomegaChanged{} class that the interval rule decides. Shares per degree of freedom, which do correct for a factor with more levels being able to spread further, keep H2 on both benchmarks, with fewer \cicids{} classes to the exporter (Table~\ref{tab:rev-robust}). The robustness analyses were added after the primary results and are not adjusted for multiplicity; we read them as checks on the verdicts, not as further tests. With three arms on UNSW-NB15 and three models, Kendall's $W$ can take only a few values (\RQfourWFlowUnsw{} at the flow unit and \RQfourWConvUnsw{} at the conversation unit on the common population), so H4 on that benchmark is weak evidence either way. The seed share of \RQtwoSeedUnswBenign{} for the UNSW-NB15 benign class is the largest on that benchmark but small next to the exporter and model shares. The RQ3 permutation test has \RQthreeClasses{} classes and little power, which is why its result is reported as supported in direction only.
